%% main.tex  –  Procedia Computer Science / ITQM conference format
%% Based on official ecrc-procs template from Elsevier

%% The '3p', 'times' and 'procedia' class options are required for Procedia CS
\documentclass[3p,times,procedia]{elsarticle}

%% The `ecrc' package must be called to make the CRC functionality available
\usepackage{ecrc}

%% ---- ecrc configuration ----
\volume{00}
\firstpage{1}
\journalname{Procedia Computer Science}
\runauth{Z. Sun, F. Shao, F. Meng}
\jid{procs}
\jnltitlelogo{Procedia Computer Science}
\CopyrightLine{2026}{Published by Elsevier Ltd.}

%% ---- packages ----
\usepackage{amssymb}
\usepackage{amsmath}
\usepackage{booktabs}
\usepackage{multirow}
\usepackage{algorithm}
\usepackage{algorithmic}
\usepackage{subcaption}
\usepackage{xcolor}
\usepackage[figuresright]{rotating}

%% ---- draft helper ----
\newcommand{\todo}[1]{\textcolor{red}{[TODO: #1]}}

\begin{document}

\begin{frontmatter}

\dochead{The 13th International Conference on Information Technology and Quantitative Management (ITQM 2026)}

\title{Credit Risk Contagion Prediction via Hierarchical Attention on
       Dynamic Heterogeneous Corporate Networks}

\author[pku]{Zhengqi Sun\fnref{fn1}}
\ead{sunzhengqi2024@stu.pku.edu.cn}

\author[pku]{Fei Shao\fnref{fn2}}
\ead{shaofei0625@stu.pku.edu.cn}

\author[pku]{Fan Meng\corref{cor1}\fnref{fn3}}
\cortext[cor1]{Corresponding author}
\ead{mengfan@pku.edu.cn}

\address[pku]{Department of Information Management, Peking University, Beijing 100871, China}

%% ---- Abstract ----
\begin{abstract}
Financial distress can propagate through inter-firm networks formed by guarantee relationships, equity linkages, and market co-movements, yet most credit risk models treat firms as independent observations.
This paper constructs a dynamic heterogeneous corporate network encompassing seven relation types for 3{,}797 Chinese A-share listed companies over 2010--2024 and investigates whether network-mediated contagion signals improve next-year distress prediction.
Statistical analysis reveals that firms in the highest neighbour-distress quintile face a 3.8$\times$ higher distress rate than those in the lowest ($\rho=0.23$, $p\approx 0$), with the effect persisting among currently profitable firms.
Logistic regressions with year fixed effects confirm that market-correlation (OR$=$1.245) and industry (OR$=$1.137) contagion channels remain significant ($p<0.001$) after controlling for firm-level financials.
We propose RA-HTGNN, a Relation-Attention Temporal Heterogeneous Graph Neural Network that combines per-relation graph convolution, relation-level attention fusion, and GRU-based temporal encoding.
While RA-HTGNN alone (AUC$=$0.795) does not surpass a Random Forest baseline with guarantee-network statistics (AUC$=$0.826), an ensemble of both achieves AUC$=$0.831, demonstrating that graph-learned representations capture complementary contagion signals beyond tabular features.
Ablation studies across edge types and GNN baselines further characterise each contagion channel's contribution.
\end{abstract}

%% ---- Keywords ----
\begin{keyword}
credit risk contagion \sep heterogeneous graph neural network \sep
temporal attention \sep guarantee network \sep corporate risk prediction
\end{keyword}

\end{frontmatter}

%% =====================================================================
\section{Introduction}
\label{sec:intro}
%% =====================================================================

Credit risk assessment has traditionally relied on firm-level financial ratios~\cite{altman1968financial, ohlson1980financial} and structural models of firm value~\cite{merton1974pricing}.
Machine learning methods---random forests, gradient boosting, and neural networks---have substantially improved prediction accuracy~\cite{barboza2017machine, moscatelli2020corporate}, yet they treat each firm as an independent observation.
In practice, financial distress propagates through inter-firm connections: a firm that guarantees a defaulting counterparty inherits contingent liabilities; companies sharing a controlling shareholder may suffer coordinated deterioration; and stocks with correlated returns reflect common exposure to systemic shocks.

Empirical studies on guarantee networks~\cite{lv2023guarantee, niu2020guarantee, ma2022credit}, supply chains~\cite{li2024credit, chen2024supply}, and equity linkages~\cite{wang2019risk} confirm that such contagion effects are economically significant.
The Chinese guarantee network, in particular, has exhibited rapid topological expansion and fragile mutual-guarantee structures that amplify systemic risk~\cite{niu2020guarantee}.
These contagion channels are heterogeneous in nature, dynamic over time, and vary in strength---properties that tabular models cannot capture.

Graph neural networks (GNNs) provide a principled framework for modelling relational dependencies~\cite{kipf2017semi, velickovic2018graph}.
Recent work has applied GNNs to guarantee-network risk prediction~\cite{cheng2020contagion, cheng2022contagion}, bankruptcy forecasting with heterogeneous relations~\cite{zhao2024combining, zheng2023qtiah, cheng2022financial}, and temporal credit modelling~\cite{wang2021temporal, chen2024temporal}.
However, most approaches either collapse all inter-firm relations into a single homogeneous graph~\cite{cheng2020contagion} or focus on a single relation type (typically guarantees).
Few works jointly model \emph{multiple heterogeneous relation types} and \emph{temporal evolution} for credit risk contagion prediction.

This paper makes three contributions:
\begin{enumerate}
\item We construct a \emph{dynamic heterogeneous corporate network} with seven relation types---guarantee, shared non-listed guarantee, equity association, equity change, co-controller, market return correlation, and industry co-membership---for 3{,}797 Chinese A-share listed companies spanning 2010--2024.
\item We provide \emph{statistical evidence of multi-channel contagion}: firms in the highest neighbour-distress quintile face a 3.8$\times$ higher distress rate, and the effect persists after year fixed effects and financial controls (market-correlation OR$=$1.245$^{***}$, industry OR$=$1.137$^{***}$), even among currently profitable firms.
\item We propose \emph{RA-HTGNN}, a hierarchical attention architecture combining per-relation graph convolution, relation-level attention fusion, and GRU temporal encoding. Its ensemble with a Random Forest achieves AUC$=$0.831, demonstrating that graph-learned contagion signals complement tabular financial features.
\end{enumerate}


%% =====================================================================
\section{Related Work}
\label{sec:related}
%% =====================================================================

\subsection{Credit Risk Prediction and Contagion}
\label{subsec:credit_risk}

Classical credit risk models rely on financial ratios~\cite{altman1968financial, ohlson1980financial} or structural models of firm value~\cite{merton1974pricing}.
Machine learning approaches---random forests, gradient boosting, and neural networks---have demonstrated consistent improvements~\cite{barboza2017machine, moscatelli2020corporate}, yet they model each firm independently.
A parallel line of research examines how risk propagates through corporate networks.
Studies on Chinese guarantee networks reveal scale-free topology and systemic fragility~\cite{niu2020guarantee, lv2023guarantee}, while epidemic-style models~\cite{ma2022credit} and supply-chain analyses~\cite{li2024credit, chen2024supply} further characterise contagion dynamics.
However, these works primarily model contagion \emph{mechanisms} rather than building end-to-end \emph{predictive} systems.

\subsection{Graph Neural Networks for Financial Risk}
\label{subsec:gnn_finance}

GNNs aggregate neighbourhood information to learn node representations~\cite{kipf2017semi, velickovic2018graph, hamilton2017inductive}.
In finance, Cheng et al.\ proposed DGANN~\cite{cheng2020contagion} for guarantee-network risk prediction and later iConReg~\cite{cheng2022contagion} for systemic crisis regulation, both operating on homogeneous guarantee graphs.
Zhao et al.~\cite{zhao2024combining} combined intra-risk and contagion-risk encoders using heterogeneous and hypergraph neural networks for bankruptcy prediction, and Zheng et al.~\cite{zheng2023qtiah} addressed class imbalance in heterogeneous GNN-based bankruptcy models.
Relational GCN~\cite{schlichtkrull2018modeling} and heterogeneous attention networks (HAN~\cite{wang2019heterogeneous}, HGT~\cite{hu2020heterogeneous}) provide architectural foundations for multi-relational graphs, but their adoption in credit risk remains limited.

\subsection{Temporal Graph Networks}
\label{subsec:temporal_gnn}

Corporate relationships evolve over time, motivating temporal graph approaches such as EvolveGCN~\cite{pareja2020evolvegcn}, DySAT~\cite{sankar2020dysat}, and TGN~\cite{rossi2020temporal}.
In finance, THGNN~\cite{xiang2022temporal} combines temporal transformers with heterogeneous graph attention for stock prediction; TemGNN~\cite{wang2021temporal} uses LSTM with interval-decayed attention for credit risk; and Zandi et al.~\cite{zandi2024attention} propose a dynamic multilayer GNN with temporal attention for loan default.
Our work differs by jointly modelling seven heterogeneous relation types with relation-level attention and GRU temporal encoding, specifically targeting multi-channel credit risk contagion in a 15-year panel of Chinese listed companies.


%% =====================================================================
\section{Problem Formulation}
\label{sec:problem}
%% =====================================================================

We model the corporate network as a dynamic heterogeneous graph $\mathcal{G} = \{G_t\}_{t=1}^{T}$, where each yearly snapshot $G_t = (\mathcal{V}, \mathcal{E}_t, \mathcal{R})$ consists of a shared node set $\mathcal{V}$ of $N$ listed companies, a time-varying edge set $\mathcal{E}_t = \bigcup_{r \in \mathcal{R}} \mathcal{E}_t^r$ comprising $|\mathcal{R}|=7$ relation types, and node feature vectors $\mathbf{x}_{v,t} \in \mathbb{R}^d$ encoding financial indicators, market information, and graph-structural statistics.

\textbf{Task.} Given a lookback window of $W$ snapshots $\{G_{t-W+1}, \ldots, G_t\}$, predict whether company $v$ will experience financial distress (defined as net profit $< 0$) in year $t+1$.


%% =====================================================================
\section{Methodology}
\label{sec:method}
%% =====================================================================

\subsection{Overview}
\label{subsec:overview}

Our framework consists of three stages: (1)~heterogeneous graph construction from multiple corporate data sources (\S\ref{subsec:graph_construction}--\ref{subsec:features}), (2)~RA-HTGNN that learns node representations through hierarchical spatial--temporal encoding (\S\ref{subsec:architecture}), and (3)~ensemble integration with a tabular Random Forest for final prediction (\S\ref{subsec:ensemble}).
The architecture draws on the hierarchical attention principle of HAN~\cite{wang2019heterogeneous}---separating node-level and semantic-level aggregation---but replaces meta-path-based aggregation with per-relation graph convolution~\cite{schlichtkrull2018modeling} and adds a GRU-based temporal encoder.

\subsection{Heterogeneous Graph Construction}
\label{subsec:graph_construction}

We construct yearly snapshots from the CSMAR database for 3{,}797 A-share listed companies.
Table~\ref{tab:edge_types} summarises the seven relation types and their characteristics.

\begin{table}[htbp]
\centering
\caption{Edge types in the heterogeneous corporate network.}\label{tab:edge_types}
\begin{tabular}{@{}llr@{}}
\toprule
Edge Type & Description & Avg. Edges/Year \\
\midrule
Guarantee & Listed-to-listed credit guarantees & $\sim$40 \\
Shared Non-listed Guar. & Co-guaranteeing a non-listed entity & $\sim$200 \\
Equity Association & Equity chain linkages & $\sim$500 \\
Equity Change & Major shareholding changes & $\sim$300 \\
Co-controller & Shared ultimate controller & $\sim$44{,}000 \\
Market Correlation & Return correlation $> \tau$ & $\sim$18{,}500 \\
Industry & Same-industry membership & $\sim$25{,}000 \\
\bottomrule
\end{tabular}
\end{table}

Guarantee and equity edges are sparse but carry direct economic meaning (contingent liability, ownership control); co-controller and market-correlation edges provide broader coverage of indirect contagion channels.
All edges are constructed independently per year from underlying CSMAR tables, forming a yearly snapshot sequence $\{G_t\}_{t=2010}^{2024}$.

\subsection{Node Feature Engineering}
\label{subsec:features}

Node features comprise: (a)~financial indicators (total assets, liabilities, revenue, cash flow, asset-liability ratio, ROA, equity ratio), (b)~current-year distress flag, (c)~credit rating information (latest rating, prospect, event count), (d)~firm attributes (market type, industry code, registered capital), (e)~graph-structural statistics (in/out degree and weighted degree per relation type), and (f)~lag-1 and year-over-year change features.
Heavy-tailed financial variables are log-transformed.
All numeric features are standardised using training-set statistics.

\subsection{RA-HTGNN Architecture}
\label{subsec:architecture}

The model consists of three hierarchical levels:

\paragraph{Level 1: Per-Relation Graph Convolution.}
Following the relation-specific parameterisation of R-GCN~\cite{schlichtkrull2018modeling}, we maintain a separate two-layer GCN~\cite{kipf2017semi} for each relation type with edge weights:
\begin{equation}
\mathbf{h}_{v,t}^{r,(l)} = \sigma\!\left(\sum_{u \in \mathcal{N}_r(v,t)} w_{uv}^r \mathbf{W}_r^{(l)} \mathbf{h}_{u,t}^{(l-1)} \right)
\end{equation}
where $\mathcal{N}_r(v,t)$ is the set of $r$-type neighbours of $v$ at time $t$, $w_{uv}^r$ is the edge weight, and $\mathbf{W}_r^{(l)}$ is a relation-specific learnable weight matrix.

\paragraph{Level 2: Relation-Level Attention.}
The outputs from $|\mathcal{R}|$ relation-specific convolutions are fused via a learnable attention mechanism:
\begin{equation}
\alpha_r = \frac{\exp(\mathbf{q}^\top \tanh(\mathbf{h}_{v,t}^r))}{\sum_{r'} \exp(\mathbf{q}^\top \tanh(\mathbf{h}_{v,t}^{r'}))}
\end{equation}
\begin{equation}
\mathbf{h}_{v,t}^{\text{spatial}} = \sum_{r \in \mathcal{R}} \alpha_r \, \mathbf{h}_{v,t}^r
\end{equation}

\paragraph{Level 3: Temporal Aggregation.}
A GRU~\cite{cho2014learning} processes the sequence of spatial representations over a window of $W$ years:
\begin{equation}
\mathbf{o}_{v,t} = \text{GRU}([\mathbf{h}_{v,t-W+1}^{\text{spatial}}, \ldots, \mathbf{h}_{v,t}^{\text{spatial}}])
\end{equation}
The final node representation concatenates three components: the last GRU hidden state (capturing sequential evolution), a temporal-attention-weighted summary (allowing the model to attend to informative historical snapshots), and a \emph{tabular residual} from an MLP applied to the raw features of the latest year (preserving firm-level signals that may be lost during graph aggregation).
A linear classifier with sigmoid activation produces the distress probability.

\subsection{Ensemble Strategy}
\label{subsec:ensemble}

The final prediction combines RA-HTGNN's predicted probability $p_{\text{GNN}}$ with a Random Forest probability $p_{\text{RF}}$:
\begin{equation}
p_{\text{ensemble}} = \alpha \cdot p_{\text{RF}} + (1-\alpha) \cdot p_{\text{GNN}}
\end{equation}
where $\alpha$ is tuned on the validation set. This leverages the RF's strength in capturing firm-level financial signals and the GNN's ability to capture relational contagion patterns.


%% =====================================================================
\section{Experiments}
\label{sec:experiments}
%% =====================================================================

\subsection{Dataset and Setup}
\label{subsec:dataset}

We use CSMAR financial data covering all A-share listed companies from 2010 to 2024.
The prediction label is next-year financial distress, defined as net profit $< 0$.
Table~\ref{tab:dataset} summarises the dataset statistics.

\begin{table}[htbp]
\centering
\caption{Dataset statistics.}\label{tab:dataset}
\begin{tabular}{@{}lr@{}}
\toprule
Item & Value \\
\midrule
Companies & 3{,}797 \\
Firm-year observations & 24{,}343 \\
Time span & 2010--2023 \\
Overall positive rate & 16.4\% \\
\midrule
Train ($\leq$2018) & 13{,}729 (pos. 9.7\%) \\
Validation (2019--2021) & 5{,}543 (pos. 18.2\%) \\
Test (2022--2024) & 5{,}071 (pos. 23.5\%) \\
\bottomrule
\end{tabular}
\end{table}

\subsection{Statistical Evidence of Contagion}
\label{subsec:stat_evidence}

Before predictive modelling, we test whether network proximity to distressed firms is associated with elevated future distress risk.

\paragraph{Chi-square Tests.}
Table~\ref{tab:contagion_chi2} reports relative risk (RR) for firms with vs.\ without a distressed neighbour under each relation type.

\begin{table}[htbp]
\centering
\caption{Association between neighbour distress and own future distress (chi-square tests).}\label{tab:contagion_chi2}
\begin{tabular}{@{}lcccc@{}}
\toprule
Edge Type & $N_{\text{connected}}$ & $P(\text{distress}|\text{dis. nbr})$ & $P(\text{distress}|\text{no dis. nbr})$ & RR \\
\midrule
Market Corr. & 22{,}909 & 0.208 & 0.090 & 2.30$^{***}$ \\
Industry & 13{,}746 & 0.169 & 0.069 & 2.44$^{***}$ \\
All Combined & 23{,}846 & 0.181 & 0.093 & 1.95$^{***}$ \\
Co-controller & 3{,}340 & 0.199 & 0.180 & 1.11 \\
Shared Non-listed Guar. & 829 & 0.214 & 0.203 & 1.06 \\
\bottomrule
\multicolumn{5}{l}{\footnotesize $^{***}$ $p<0.001$} \\
\end{tabular}
\end{table}

\paragraph{Dose--Response Relationship.}
We divide firms into quintiles by aggregate neighbour distress ratio.
The highest quintile exhibits a distress rate of 32.5\%, compared with 8.6\% for the lowest---a 3.8$\times$ difference (Spearman $\rho = 0.230$, $p \approx 0$).
Critically, this gradient persists among currently profitable firms (RR$=$1.63--1.96, $p<0.001$), demonstrating early-warning value.

\paragraph{Logistic Regression with Year Fixed Effects.}
To control for firm-level financials and time trends, we estimate logistic regressions on the full sample ($N=$24{,}343) with year dummies.
Table~\ref{tab:logistic} reports contagion variable odds ratios from the multivariate model.

\begin{table}[htbp]
\centering
\caption{Logistic regression with year fixed effects: odds ratios for lagged contagion variables (Model~B).}\label{tab:logistic}
\begin{tabular}{@{}lcccc@{}}
\toprule
Contagion Channel & OR & 95\% CI & $z$ & $p$ \\
\midrule
Market Correlation & 1.245$^{***}$ & [1.201, 1.291] & 11.85 & $<$0.001 \\
Industry & 1.137$^{***}$ & [1.097, 1.178] & 7.05 & $<$0.001 \\
Shared Non-listed Guar. & 1.031 & [0.999, 1.064] & 1.88 & 0.061 \\
Equity Change & 1.000 & [0.965, 1.037] & 0.01 & 0.993 \\
Equity Association & 0.986 & [0.965, 1.008] & $-$1.23 & 0.217 \\
Co-controller & 1.013 & [0.983, 1.044] & 0.84 & 0.399 \\
\bottomrule
\multicolumn{5}{l}{\footnotesize Controls: net profit, ROA, equity ratio, cashflow/assets, asset-liability ratio,} \\
\multicolumn{5}{l}{\footnotesize \phantom{Controls: }current loss flag, log(total assets), year dummies. Pseudo-$R^2=0.196$.} \\
\multicolumn{5}{l}{\footnotesize LR test vs.\ financial-only model: $\chi^2=226.7$, $p=3.9\times10^{-46}$.} \\
\multicolumn{5}{l}{\footnotesize $^{***}$ $p<0.001$. Guarantee excluded due to zero variance.} \\
\end{tabular}
\end{table}

Among currently profitable firms (Model~D, $N=$20{,}345), market correlation (OR$=$1.296$^{***}$) and industry (OR$=$1.135$^{***}$) remain significant, and shared non-listed guarantees become significant (OR$=$1.049$^{**}$, $p=0.006$).


\subsection{Baselines}
\label{subsec:baselines}

We compare the following models:
\begin{itemize}
\item \textbf{Logistic Regression} and \textbf{Random Forest} (RF): tabular models trained on financial features, with an extended variant (L1) that includes guarantee-network graph statistics.
\item \textbf{MLP}: a feedforward network using the same node features as RA-HTGNN but without graph message passing, serving as an ablation of the graph component.
\item \textbf{HomoGCN} and \textbf{HomoGAT}: standard two-layer GCN~\cite{kipf2017semi} and GAT~\cite{velickovic2018graph} on the union graph (all edge types merged into a single relation), combined with a GRU temporal encoder.
\item \textbf{RA-HTGNN} (ours): per-relation graph convolution with relation-level attention and GRU temporal encoding.
\end{itemize}

\subsection{Implementation Details}
\label{subsec:impl}

All GNN models use hidden dimension 96, dropout 0.25, temporal window $W=3$, AdamW~\cite{loshchilov2019decoupled} with learning rate $8\times10^{-4}$ and weight decay $5\times10^{-4}$, BCE loss with positive class weight, and early stopping with patience 20 over 120 epochs.
RF uses 500 trees with balanced class weights and max depth 12.
Experiments are implemented using DGL~\cite{wang2019dgl} and PyTorch.


\subsection{Main Results}
\label{subsec:main_results}

Table~\ref{tab:main_results} compares all models on the test set.

\begin{table}[htbp]
\centering
\caption{Main prediction results on the test set (2022--2024).}\label{tab:main_results}
\begin{tabular}{@{}llcc@{}}
\toprule
Model & Features & Test AUC & Test AP \\
\midrule
Logistic Regression & Financial (L0) & 0.769 & 0.540 \\
Random Forest & Financial (L0) & 0.821 & 0.618 \\
Random Forest & Financial + Guarantee (L1) & 0.826 & 0.610 \\
\midrule
MLP (no graph) & All node features & 0.797 & 0.603 \\
HomoGCN & All edges merged & 0.795 & 0.586 \\
HomoGAT & All edges merged & 0.796 & 0.576 \\
RA-HTGNN (ours) & 6 edge types & 0.795 & 0.578 \\
\midrule
\textbf{RF (L1) + RA-HTGNN Ensemble} & $\alpha=0.7$ & \textbf{0.831} & \textbf{0.623} \\
\bottomrule
\end{tabular}
\end{table}

Three key observations emerge.
First, the RF baseline is a strong predictor: even without graph information (L0, AUC$=$0.821), it outperforms all neural models, and adding guarantee-network statistics (L1, AUC$=$0.826) yields the best single-model result, consistent with prior findings on the dominance of accounting-based indicators~\cite{barboza2017machine, duan2012multiperiod}.
Second, all neural models---MLP, HomoGCN, HomoGAT, and RA-HTGNN---achieve similar test AUC ($\sim$0.795--0.797), below the RF baseline. This gap likely reflects the moderate sample size ($\sim$24K firm-years) and temporal distribution shift (positive rate rising from 9.7\% to 23.5\%), which limit end-to-end neural training.
Third, the ensemble of the L1 RF and RA-HTGNN yields the best overall performance (AUC$=$0.831), improving upon the strongest standalone model by $+$0.5 percentage points and confirming that the GNN captures complementary contagion signals not present in tabular features.
Notably, the homogeneous GNN baselines do not outperform the MLP, suggesting that na\"{i}vely merging heterogeneous edges introduces noise rather than useful relational structure.


\subsection{Ablation Study}
\label{subsec:ablation}

\paragraph{Edge Type Ablation.}
Table~\ref{tab:ablation} reports RA-HTGNN performance with different relation subsets.

\begin{table}[htbp]
\centering
\caption{Ablation: RA-HTGNN test AUC by edge subset and corresponding ensemble (seed 42).}\label{tab:ablation}
\begin{tabular}{@{}lcc@{}}
\toprule
Relations Used & Test AUC & Ensemble AUC \\
\midrule
All 7 types & 0.790 & 0.828 \\
6 types (no industry) & 0.795 & \textbf{0.831} \\
Guarantee + Shared guar. & 0.808 & 0.830 \\
Guarantee only & 0.808 & 0.830 \\
Equity assoc. + Equity change & 0.804 & 0.829 \\
Market correlation only & 0.799 & 0.830 \\
5 types (no market corr.) & 0.801 & 0.829 \\
\bottomrule
\end{tabular}
\end{table}

A counter-intuitive finding is that models with fewer edge types achieve higher standalone AUC.
The ``guarantee only'' model (AUC$=$0.808) outperforms the 6-edge variant (0.795) despite having only $\sim$40 edges per year.
This occurs because RA-HTGNN includes a tabular-residual MLP path: with very sparse edges, the GCN branch contributes minimal aggregation noise, so the model effectively reduces to an MLP with a small graph bonus.
Adding dense, low-homophily edges (market correlation, co-controller) introduces neighbourhood noise that degrades the spatial representations.
However, ensemble performance is remarkably stable across all configurations (AUC 0.828--0.831), indicating that even noisy graph representations provide complementary signal when combined with the RF.

\paragraph{Multi-Seed Stability.}
Over 4 seeds (42, 123, 7, 1), RA-HTGNN with the best 6-edge configuration achieves test AUC $0.793 \pm 0.005$ and ensemble AUC $0.829 \pm 0.002$, confirming that the ensemble improvement is consistent and not an artefact of seed selection.


%% =====================================================================
\section{Analysis and Discussion}
\label{sec:analysis}
%% =====================================================================

\subsection{Why Contagion is Statistically Significant but Prediction-Limited}
\label{subsec:why_limited}

The apparent paradox---strong statistical contagion effects (RR up to 2.44, OR up to 1.296) yet limited standalone GNN gain---is explained by the causal chain:
\emph{contagion risk exposure} $\rightarrow$ \emph{financial deterioration} $\rightarrow$ \emph{observed financial indicators}.
Since the RF already uses downstream financial indicators (net profit, ROA, asset-liability ratio) as features, contagion information is largely subsumed by these endogenous variables.
The residual contagion signal that RA-HTGNN captures---network-structural patterns not yet reflected in accounting numbers---enables the ensemble to improve, but the marginal gain is bounded by this information overlap.
This finding echoes the broader observation that network features tend to provide incremental rather than transformative improvements when strong firm-level predictors are available~\cite{wang2022credit}.

\subsection{Interpretation of Contagion Channels}
\label{subsec:channels}

Market return correlation (OR$=$1.245) likely captures common exposure to macroeconomic and sector-level shocks, consistent with the systemic risk literature~\cite{niu2020guarantee}.
Industry co-membership (OR$=$1.137) reflects demand-side contagion within product markets, where distress in one firm signals deteriorating industry conditions.
Guarantee linkages, while economically the most direct contagion channel, are too sparse among listed companies ($\sim$40 edges/year) to yield robust statistical signals---a limitation also noted in studies restricted to listed-firm networks~\cite{lv2023guarantee}.
Paradoxically, this sparsity benefits the GNN: the ablation shows that guarantee-only models achieve higher standalone AUC (0.808) than multi-edge variants (0.795), because sparse edges avoid the aggregation noise introduced by dense, low-homophily connections.
The ensemble, however, is robust across all edge configurations (AUC 0.828--0.831), suggesting that the complementary signal extracted by graph learning is more about structural diversity than any single edge type.

\subsection{Limitations}
\label{subsec:limitations}

Several limitations should be noted.
(1)~Our graph only covers listed companies; incorporating non-listed entities connected through guarantee chains~\cite{cheng2022contagion} could substantially strengthen contagion signals.
(2)~Market return correlation may reflect common factor exposure rather than directional risk transmission; disentangling these requires causal identification strategies beyond the scope of this study.
(3)~The positive rate increases from 9.7\% (train) to 23.5\% (test), reflecting macroeconomic deterioration during 2022--2024, which introduces temporal distribution shift.
(4)~Yearly snapshots cannot capture within-year contagion dynamics; continuous-time approaches~\cite{rossi2020temporal} may better model rapid risk propagation.


%% =====================================================================
\section{Conclusion}
\label{sec:conclusion}
%% =====================================================================

This paper investigated credit risk contagion prediction through dynamic heterogeneous corporate networks covering seven relation types for 3{,}797 Chinese A-share companies over 2010--2024.
We established statistically significant contagion effects across multiple channels, with market-correlation and industry co-membership contributing most strongly after controlling for firm-level financials.
Our proposed RA-HTGNN, combined with a Random Forest in an ensemble, achieves AUC$=$0.831, demonstrating that graph-learned contagion representations provide complementary predictive value beyond tabular financial features.
The finding that contagion signals are largely mediated through observable financial indicators has practical implications: network monitoring can serve as an early-warning system even when its direct predictive margin is modest.

Future work will pursue three directions: (1)~incorporating non-listed entities into the guarantee network to recover stronger contagion signals~\cite{cheng2022contagion}, (2)~adopting continuous-time temporal graph architectures~\cite{rossi2020temporal} for finer-grained dynamics, and (3)~exploring attention-based explainability to identify the most influential contagion paths for regulatory decision-making.


%% =====================================================================
%% References
%% =====================================================================
\bibliographystyle{elsarticle-num}
\bibliography{references}

\end{document}
